\chapter{設計・計画}
\section{構造設計は力の流れと性能を結び付ける}
橋、機械の軸、容器や支持柱は、与えられた力を接合部から支持部へ伝える。
設計では材質を強くするだけでなく、荷重がどう流れ、変形がどこに集まり、支持機構がどの条件で失われるかを調べる。
同じ荷重でも細い柱は材料の破断前に座屈し、強い梁でも過大なたわみで使えなくなる。
本節群は構造強度の初稿である。他の設計・計画領域は後続の原典確認で加筆する。

\section{強度・剛性・安定を分ける}
\label{theory:structural:modes}
強度は破損・降伏等に耐える能力、剛性は荷重に対する変形の小ささ、安定は小さな擾乱が増幅せず平衡を維持する性質である。
ヤング率 \(E\) が同じでも降伏強度は異なり得る。逆に強度の高い材料が必ず高剛性であるとはいえない。
静的な荷重支持機構の喪失、圧縮による座屈、周期外力による共振を区別する。
共振の比較対象は外力周波数と固有振動数であり、応力と振動数を大小比較することは次元的にも誤る。
短く太い柱では圧縮降伏・圧壊が支配し得るが、引張破断強度をそのまま圧縮破損の閾値としない。

\section{引張試験の測定値から定義を作る}
\label{theory:structural:tensile}
試験機が直接取得するのは荷重 \(F\) と標点間の伸び \(\Delta L\) である。
初期断面積 \(A_0\)、初期標点距離 \(L_0\) を用いて
\[
 \sigma_{\mathrm n}=\frac{F}{A_0},\qquad
 \varepsilon_{\mathrm n}=\frac{\Delta L}{L_0}
\]
と規格化すると、試験片の初期寸法の影響を分けて比較できる。
除荷で回復する変形が弾性、残る変形が塑性である。弾性という語だけでは線形比例まで保証しない。
線形範囲では \(\sigma=E\varepsilon\) である\cite{struct_mlit_glossary}。
延性は破断まで塑性変形を許す能力であり、塑性と同じ語ではない。
単調引張の最大公称応力
\[
 \sigma_{\mathrm B}=\frac{\max F}{A_0}
\]
が引張強さである。くびれ後の瞬間断面積 \(A(t)\) を使う真応力 \(F/A(t)\)、破断時の公称応力、
繰返し負荷での疲労限度を混同しない。平成29年度の「応力」の語には、初期断面積を用いる前提を明示して読む。

\section{許容応力と安全率}
\label{theory:structural:allowable}
材料の限界応力を \(\sigma_{\mathrm{str}}\)、安全率を \(n>1\) とした単純なモデルでは
\[
 \sigma_{\mathrm{allow}}=\frac{\sigma_{\mathrm{str}}}{n},
 \qquad \sigma_{\mathrm{working}}\le \sigma_{\mathrm{allow}}
\]
である。安全側へ余裕を取るため許容値は限界値より小さくなる。
ヤング率も応力単位を持つため、強度をヤング率で割った比を許容応力と呼ぶことはできない。
実設計では何を限界応力にするか、温度、荷重種類、疲労、座屈等を分ける必要がある。
安全率を掛けて許容値を大きくする誤りと、材料強度だけで全ての破損形態を代表させる誤りを避ける。

\section{限界状態は一つではない}
\label{theory:structural:limit}
限界状態は要求された性能を満たせなくなる境界である。
通常の使用に必要な変位・振動等を扱う使用限界、最大耐荷性能を扱う終局限界、
繰返し負荷の損傷を扱う疲労限界を分ける。
破壊していなくても利用できなくなる状態があるため、終局だけを照査しても十分ではない。
国土交通省の空港構造設計要領は修復限界も扱い、疲労を終局に含める場合があることを明記しつつ別に分類する\cite{struct_mlit_limit}。
これはすべての構造種別で分類名と係数が同一という意味ではない。過去問で指定された対応と規準の粒度を保存する。

\section{不静定構造の余力とその限界}
\label{theory:structural:redundancy}
静定構造の反力・断面力は適切な平衡式から求められる。不静定構造では余剰な未知力があり、
変形の適合条件と材料の構成関係も用いる。
延性的な部材が局所降伏した後、他の経路へ荷重を再配分できると構造全体の余力が生まれる。
この説明には接合部の能力、十分な変形能力、荷重経路の存在が必要である。
塑性化による材料非線形、たわみによる幾何非線形を扱わず、常に線形比例を仮定して損傷後を解析することはできない。
\begin{tcolorbox}[title=部分損傷から全体安全へ飛躍しない]
不静定次数が高いことだけでは安全性の保証にならない。少数の塑性ヒンジによる機構形成、
全体座屈や接合部破壊で支持能力が失われる場合がある。
平成24年度の初回・公式5は保持するが、「すべての部材が破壊するまで安全」という一般命題へ拡張しない。
\end{tcolorbox}

\section{梁の曲げ：断面二次モーメントと断面係数}
\label{theory:structural:beam}
細長い均質梁の小変形・線形弾性では、平面断面が変形後も平面を保つと近似する。
中立軸からの距離を \(y\)、曲率を \(\kappa\) とすると軸ひずみは \(-y\kappa\)、応力は \(-Ey\kappa\) となる。
断面にわたりモーメントを積分して
\[
 I=\int_A y^2\,dA,\qquad |M|=EI|\kappa|,\qquad
 \sigma_{\max}=\frac{|M|c}{I}=\frac{|M|}{Z},\qquad Z=\frac{I}{c}
\]
を得る。\(c\) は外縁までの距離である。\(I\) の単位は長さの4乗、\(Z\) は3乗である。
断面を中立軸から離して配置すると曲げへの抵抗が増す。これは材料そのものの強度を変更する操作ではない。
梁理論の仮定、断面力、境界条件は大学の公開講義と対応する\cite{struct_mit_beam}。

\section{単純支持梁の等分布荷重を解く}
一定の荷重密度 \(q\)、長さ \(L\) の梁の両端を単純支持する。対称性と平衡から反力は両端 \(qL/2\) である。
左端からの位置 \(x\) でのモーメントは
\[
 M(x)=\frac q2 x(L-x),\qquad M_{\max}=\frac{qL^2}{8}.
\]
下向きたわみを正とする符号規約で \(EI v''=-M\) とし、\(v(0)=v(L)=0\) を用いて積分すると
\[
 v(x)=\frac{q}{24EI}\bigl(x^4-2Lx^3+L^3x\bigr),\qquad
 v_{\max}=v(L/2)=\frac{5qL^4}{384EI}
\]
を得る。最大応力は材料と断面係数、最大たわみは \(EI\) と荷重・長さに依存する。
大たわみ、短く太い梁のせん断変形、断面の局所座屈をこの式に含めたとはしない。

\section{寸法のべき則は固定条件を含む}
\label{theory:structural:scaling}
円断面の直径 \(d\)、矩形断面の幅 \(b\) と荷重方向の高さ \(h\) に対して
\[
 I_{\mathrm{circle}}=\frac{\pi d^4}{64},\quad
 Z_{\mathrm{circle}}=\frac{\pi d^3}{32},\qquad
 I_{\mathrm{rectangle}}=\frac{bh^3}{12},\quad
 Z_{\mathrm{rectangle}}=\frac{bh^2}{6}.
\]
\(q,L,E\) を固定して円梁の径だけを変えると \(\sigma_{\max}\propto d^{-3}\)、\(v_{\max}\propto d^{-4}\) である。
矩形の幅も固定すれば \(\sigma_{\max}\propto h^{-2}\) となる。断面二次モーメントをそのまま応力の分母へ入れると、
外縁距離の1乗を見落とす。
\begin{tcolorbox}[title=自重を固定荷重として扱わない]
自重だけなら \(q=\rho Ag\) である\cite{struct_mit_beam}。円梁の \(A\propto d^2\) も同時に変わるため、
同じ \(\rho,L,E\) では自重応力は \(d^{-1}\)、たわみは \(d^{-2}\) となる。
令和4・6年度の解は通常の固定 \(q\) 比較として示し、寸法を変えれば全ての荷重が一定であるとは主張しない。
\end{tcolorbox}

\section{座屈は圧縮平衡の安定性喪失である}
\label{theory:structural:buckling}
引張は細い定規を直線へ戻す方向に働くが、圧縮は小さな横たわみを増幅し得る。
理想的な等断面柱で横たわみを \(v(x)\)、圧縮力を \(P\) とすると、両端ピンの微小たわみモデルでは
\[
 EI v''+Pv=0,\qquad v(0)=v(L)=0
\]
となる。非自明解は \(v\propto\sin(m\pi x/L)\) であり、最小の固有値は
\[
 P_{\mathrm{cr}}=\frac{\pi^2EI}{(KL)^2}
\]
である。\(K\) は端部拘束に応じた有効長係数である。両端ピン1、固定・自由2、両端固定0.5という理想化を用いる。
同じ拘束と断面なら長さの逆2乗、円断面なら直径の4乗に依存する。
弾性式に降伏強度は直接出ないが、臨界応力 \(P_{\mathrm{cr}}/A\) が弾性域にあることが適用条件である。
細長比が小さい領域の降伏・圧壊をオイラー式だけで判定しない\cite{struct_mit_buckling}。
初期曲がりや荷重偏心のある実柱には、理想臨界前からたわみがある。図の突然の横変形を共振や剥離と呼ばない。

\section{疲労は振幅と平均応力を分ける}
\label{theory:structural:fatigue}
繰返しの荷重・ひずみが損傷を蓄積し、静的破壊強度より小さい負荷でも破損することがある。
応力の最大・最小から
\[
 \sigma_a=\frac{\sigma_{\max}-\sigma_{\min}}{2},\qquad
 \sigma_m=\frac{\sigma_{\max}+\sigma_{\min}}{2}
\]
を定義する。振幅と平均を混ぜず、同じ振幅の比較では通常の金属疲労モデルで平均引張応力が不利となる。
修正Goodman線はその影響を扱う設計近似であるが、材料・残留応力・塑性・多軸条件を問わず同一の寿命式にはならない。
データの温度、応力比、ひずみ制御か荷重制御か、対象繰返し数を読む\cite{struct_nims_data,struct_nasa_mean}。
疲労限度が存在するという理想化もすべての金属や超高サイクル域へ一般化しない。

\section{クリープは時間と温度の条件を持つ}
\label{theory:structural:creep}
金属材料が高温で荷重を受け続けると、一定の負荷でも塑性変形が時間とともに進む。
瞬間的なヤング率や単調引張強度だけでは長期の変形・破断時間を決められない。
NIMSの資料は高温機器の設計に、温度・荷重・経過時間を揃えた長時間試験が必要なことを示す\cite{struct_nims_data}。
高強度という語を高温という条件へ置換しない。また低温から常に全材料の時間依存変形がゼロとなるという定理にはしない。

\section{応力集中は局所と公称の差である}
\label{theory:structural:concentration}
穴、切欠き、段差等の幾何変化は応力の流れを曲げる。局所最大と遠方の公称応力の比で集中を評価する。
半径を緩やかにする等の設計で集中を減らすことがあるが、その値は形状・荷重・境界条件に依存する。
一軸応力ではvon Mises相当応力は軸応力の絶対値に一致するが、一般の多軸場は
\[
 \sigma_{\mathrm{VM}}=
 \sqrt{\frac{(\sigma_1-\sigma_2)^2+(\sigma_2-\sigma_3)^2+(\sigma_3-\sigma_1)^2}{2}}
\]
で評価する。主応力 \(\sigma_i\) の静水圧部分はこの値へ寄与しない。
断面を急変させたという言葉だけで、どの荷重・どの点の相当応力も必ず高くなるとは断言しない。

\section{比強度と価格当たり性能}
\label{theory:structural:specific}
軽量な荷重支持を考えるとき、比強度 \(\sigma/\rho\) は同じ質量でどれだけ引張荷重を支えられるかの指標となる。
価格が質量当たり \(p\) なら、単位比強度当たりの価格は \(p/(\sigma/\rho)\) である。
二つの材料で比強度が何倍かを求めてから、価格比をその倍率で割る。
令和5年度の値ではCFRPの比強度が鋼の24.6875倍、価格が266.6667倍なので、
単位比強度当たりの価格は約10.80倍となる。強度優位と費用優位は別である。
この価格は設問所与であり現在の市場価格ではない。繊維方向、圧縮・接合、加工、疲労等の性能をこの一指標へ圧縮しない。

\section{強度と荷重の確率には同時分布が必要である}
\label{theory:structural:reliability}
強度 \(R\) と作用 \(S\) にばらつきがあるモデルでは、破壊事象を \(S>R\) とする。
一般に
\[
 P_{\mathrm f}=\iint_{s>r} f_{R,S}(r,s)\,ds\,dr
\]
であり、独立なら \(f_{R,S}=f_R f_S\) と因数分解できる。
平成29年度の \(R\sim U[18,23]\)、\(S\sim U[10,20]\) について独立を補うと、
破壊領域は一辺2の直角三角形、支持矩形の面積は50なので \(P_{\mathrm f}=2/50=0.04\) となる。
\begin{tcolorbox}[title=周辺分布だけでは決まらない]
原文には独立性の明記がない。同じ \(U\sim U[0,1]\) から
\(R=18+5U,\ S=10+10U\) とすると破壊確率0である。
一方 \(R=18+5U,\ S=20-10U\) とすると破壊条件は \(U<2/15\) で、確率2/15となる。
どちらも所与の周辺一様分布を持つ。公式0.04を保持しつつ、独立モデルを用いた条件付き解であることを残す。
\end{tcolorbox}

\section{安全率の費用最小化と規範的判断を分ける}
\label{theory:structural:optimization}
対策を増やすと直接費用は増し、事故等の損害費用は減るという所与のモデルを考える。
\[
 C(x)=\frac{a}{1+bx}+cx,\quad
 C'(x)=-\frac{ab}{(1+bx)^2}+c,\quad
 C''(x)=\frac{2ab^2}{(1+bx)^3}.
\]
\(a,b,c>0\) かつ \(1+bx>0\) の域では凸であり、停留点が定義域内なら一意の最小である。
根が域外なら許容端点を比較し、開区間では端点に最小が実現しない場合もある。
令和元年度再試験の \((a,b,c)=(90,4,10)\)、\(x\ge1\) では \(x=1.25\)、
令和4年度の \((9,1,1)\)、\(x>1\) では \(x=2\) である。全費用の単位を億円にそろえる。
この数学的最小化は与式の範囲内の答えであり、法的な最低要求、重大危害、ALARPの著しい不均衡を
単純な費用関数だけで判断できるという主張ではない。

\section{典型的な問われ方と確認順序}
語の穴埋めでは測定値と規格化した量、材料物性と形状の効果、荷重形態と破損形態を分ける。
寸法の比例問題では断面二次モーメントと断面係数を先に書き、固定した荷重を確認する。
最適化では定義域と単位、確率では独立性と領域を先に読む。
同じ定規や柱の題材でも年度によって候補の順が違うため、正答番号だけを移さない。
個別問題から以上の共通理論へ参照し、全問の原典確認と分類を終えるまでは頻度や将来の出題を断定しない。

\section{信頼性を確率モデルとして設計する背景}
\label{theory:dependability:background}
強度を満たしても、制御、電源、通信など一つの必要機能が止まれば装置全体が使えなくなる。
一方、代替経路を設ければ一部の停止を許せる。この違いを、各部の成功又は故障という事象の論理関係へ写すのが信頼性設計である。
対象機能、環境条件、観測期間を先に定め、その期間内に要求機能を維持する確率を信頼度とする。
部品の配置図は電気回路そのものとは限らず、機能が成立する条件を表す。

\section{信頼度、寿命分布、故障率}
\label{theory:dependability:definition}
最初の故障までの寿命を非負確率変数\(T\)とし、
\[
 R(t)=P(T>t),\qquad F(t)=1-R(t),\qquad
 h(t)=\frac{f(t)}{R(t)}
\]
と定める。最後の式は密度\(f\)が存在し\(R(t)>0\)の範囲での条件付き故障率である。
信頼度は無次元、故障率は時間の逆数の次元を持つ。一定故障率\(\lambda\)を仮定した場合に限り\(R(t)=e^{-\lambda t}\)となる。
直並列の論理を扱う設問は時点又は期間を固定した信頼度を与えるため、指数寿命分布を無断で追加する必要はない。
\cite{reliability_nist_series}

\section{直列系：全ての必要機能が成功する}
\label{theory:dependability:series}
要素\(i\)が成功する事象を\(S_i\)とする。全てを必要とする系の成功は\(\bigcap_i S_i\)である。
一般には条件付き確率の積であり、独立を仮定したときだけ
\[
 R_{\mathrm{series}}=P\left(\bigcap_i S_i\right)=\prod_i r_i,
 \qquad r_i=P(S_i)
\]
となる。同じ信頼度\(r\)の要素を増やすと\(r^n\)であり、\(0<r<1\)では低下する。
最初の一要素故障が系を止める非修復モデルの条件である。\cite{reliability_nist_series}

\section{並列系：少なくとも一つの経路が成功する}
\label{theory:dependability:parallel}
系の停止は全経路の停止である。補集合を使うと、独立な要素について
\[
 R_{\mathrm{parallel}}
 =1-P\left(\bigcap_i S_i^{\mathrm c}\right)
 =1-\prod_i(1-r_i)
\]
となる。同一要素\(n\)個なら\(1-(1-r)^n\)である。
「並列」という語だけで成功率を足すと、二経路が同時成功する重複を数えて1を超え得る。
待機予備の切替故障、共通電源、負荷再配分による故障率変化は、全要素が独立して動くこの模型から分ける。\cite{reliability_nist_parallel}

\section{独立性がない場合と共通原因故障}
\label{theory:dependability:independence}
二要素では、周辺成功率\(p,q\)が分かっても
\[
 \max(0,p+q-1)\le P(S_1\cap S_2)\le\min(p,q)
\]
の幅が残る。下限は和事象の確率が1以下、上限は共通事象が各事象に含まれることから出る。
同じ地震、冷却系、電源又はソフトウェアの誤りで停止すれば、機器の台数が二つでも独立とは限らない。
例えば信頼度0.9の二要素が完全に同じ成功事象を持つなら並列系も0.9であり、独立模型の0.99にはならない。
原典が独立を明示しない計算問題では、公式の想定模型と条件不足を併記する。

\section{混合系の縮約と未知信頼度の逆算}
\label{theory:dependability:mixed}
独立な別要素からなる部分系を内側から一つの成功事象へまとめ、次の接続を読む。
縮約後の部分系同士が要素を共有すると、元の部品が独立でも部分系は独立ではない。
未知信頼度\(x\)を逆算する場合はまず\(0\le x\le1\)を課し、系の信頼度との区別を保つ。
二系の等式から\(x\)を割る操作は\(x>0\)が確認された場合に限る。
原典に非零条件がないときは、零根を都合よく消さない。全体の信頼度を問う候補へ単体\(x\)を入れない。

\section{要求信頼度を満たす最小冗長数}
\label{theory:dependability:redundancy}
独立並列の要素信頼度\(r\)、その外に直列で必要な一要素の信頼度\(s\)を置けば
\[
 R_n=s\{1-(1-r)^n\}.
\]
要求\(R_n\ge R_*\)を変形し、\(0<r<1\)かつ\(R_*<s\)なら
\[
 (1-r)^n\le1-\frac{R_*}{s},\qquad
 n\ge\frac{\log(1-R_*/s)}{\log(1-r)}.
\]
負の分母で割るため不等号の向きに注意する。最小整数の前後を直接代入して検算する。
有限の\(n\)で\(R_n<s\)なので、外の直列要素が達成可能な上限を作る。
これを改善せず並列部だけ増やしても要求を満たせない場合がある。

\section{多数決冗長系}
\label{theory:dependability:voting}
\(n\)個のうち任意の\(k\)個以上が成功すればよい場合、独立・同一信頼度\(r\)の模型では
\[
 R_{k/n}=\sum_{j=k}^{n}\binom{n}{j}r^j(1-r)^{n-j}.
\]
二項係数は、成功する要素の組合せを数えるものである。成功個数が異なる事象は排反なので和を取れる。
2/3系は\(3r^2(1-r)+r^3\)であり、全並列の1/3系とも全直列の3/3系とも異なる。
判定器自身、同じ故障出力を返す共通原因、異なる単体信頼度は別の模型を要する。\cite{reliability_nist_voting}

\section{故障の木：原因事象から上位事象へ}
\label{theory:dependability:fta}
故障の木解析（FTA）は、望ましくない上位事象がどの下位事象の論理結合で起きるかを表す。
独立な下位故障確率を\(p_i\)とした場合、ANDは\(\prod_i p_i\)、ORは\(1-\prod_i(1-p_i)\)となる。
信頼性ブロック図の成功事象とFTAの故障事象を区別する。形が似ていても、どちらの事象を入力としているかを先に読む。
末端から段階的に縮約し、同じ原因が複数枝へ再出現する場合は事象を重複した独立変数にしない。
独立でない二事象のORは常に\(p_1+p_2-P(E_1\cap E_2)\)であり、最後の項を勝手に\(p_1p_2\)へ置換できない。

\section{修復可能な系の稼働率}
\label{theory:dependability:availability}
信頼度は途中で故障しない確率であるが、稼働率は故障後の修復も含めて所要時点又は期間に使える割合である。
平均運転時間を\(m_{\mathrm u}\)、平均修復を含む非稼働時間を\(m_{\mathrm d}\)とする。
交互の運転・停止周期が安定し有限平均を持つ模型で、長期の時間割合は
\[
 \mathcal A=\frac{m_{\mathrm u}}{m_{\mathrm u}+m_{\mathrm d}}.
\]
多周期の運転時間総和を周期時間総和で割ると、平均の比としてこの式が出る。指数分布は必要条件ではない。
試験でMTBFを平均運転時間、MTTRを平均修復時間とする場合はMTBF/(MTBF+MTTR)となる。
同じ名称でも暦時間の故障間隔を含める定義を代入すると分母を二重計上する。
物流待ち、点検、計画停止を含む運用稼働率には、その時間範囲に対応する分母を使う。

\section{保全：故障前後と処置時期の判断基準}
\label{theory:dependability:maintenance}
予防保全は故障を避ける目的で故障前に処置し、事後保全は故障後に機能を回復する。
時間計画保全は暦又は運転時間等から時期を予定し、状態監視保全は測定や点検で把握する劣化状態から処置を決める。
点検そのものが定期的でも、処置の判断が状態に基づくなら時間計画の交換とは区別される。
国交省の設備保全資料は、重大な機能影響と状態監視の可能性で三種の使い分けを示す。\cite{maintenance_mlit_2013}
全設備へ同じ交換周期を当てはめず、故障の帰結、兆候の観測可能性、保全の費用を検討する。
JIS・IECの版ごとの全用語一致はここでは未確認である。

\section{安全側、縮退、誤操作対策、耐故障性}
\label{theory:dependability:handling}
故障時に安全を維持するフェールセーフ、影響を受ける機能を優先順位付きで終了するフェールソフト、
人為的な誤りを重大な失敗へ結び付けないフールプルーフ、想定する故障があっても機能するフォールトトレランスを分ける。
例えば機械の不意の再起動を防ぐ仕組みは故障時の安全条件を定め、開いたガードで起動しない仕組みは誤操作への備えにもなる。\cite{design_mhlw_failsafe,design_mhlw_foolproof}
耐故障性には検出・隔離・回復や冗長性が関わり、影響を外へ見せないフォールトマスキングはその一方式である。
航空宇宙のソフトウェアでも故障の予防と発生後の影響緩和を分けて設計する。\cite{design_nasa_tolerance}
安全に停止できる設備と機能継続が安全上必要な設備では要求が異なるため、「止めれば常に安全」とはしない。

\section{安全率の値を比べる前に定義を比べる}
\label{theory:dependability:factor}
安全率\(n\)は、基準強さ\(\sigma_{\mathrm{ref}}\)から許容応力を
\(\sigma_{\mathrm{allow}}=\sigma_{\mathrm{ref}}/n\)とする模型で、無次元の余裕を与える。
基準が引張強さか降伏強さか疲労限度かを先に指定する。同じ材料でも荷重履歴や破損形態が違えば値を転用しない。
設問の2以上又は1.3～2.0という代表値だけで、全設計に共通する法定下限を作ることはできない。
\begin{tcolorbox}[title=対象と資料時点を限定した例]
2011年の米国航空機規則25.303では、外力としての限界荷重に原則1.5を掛ける。ロケット全般の規定ではない。
厚労省のクレーン等安全規則213条では、玉掛けワイヤロープの切断荷重/最大荷重を6以上とする。
食品安全委員会の2008年ADI説明は、残留農薬等の無毒性量を通常100で割る。治療薬全般の一律係数ではない。
国交省の平成25年改正反映告示表では、対象条件を満たすエレベーター主索の常時使用時4.0、設置時5.0であり端部の値は別である。
この告示の2011当時版及び現行までの全改正履歴は未照合である。
\end{tcolorbox}
\cite{design_faa_2011,design_mhlw_crane,design_fsc_adi,design_mlit_elevator}
異なる分母と対象を持つ係数の大小を、危険又は安全の単純順位へ変換しない。

\section{設計管理と製造物責任法の安全性}
\label{theory:dependability:liability}
大量生産品では、利用者が製造過程の過失を立証することが難しい。この背景で製造物の欠陥に基づく責任を定め、被害者の保護を図るのが製造物責任法の目的である。
第2条の欠陥は、特性、通常予見される使用形態、引渡時期などを考慮した通常有すべき安全性の不足を意味する。
機能性の不足、故障、破損だけで直ちに全て法定の欠陥になるわけではなく、絶対安全の要求でもない。
消費者庁は製造上、設計上、指示・警告上の三側面を説明している。\cite{design_caa_pl1,design_caa_pl2}
設計者は必要機能だけでなく、合理的に予見する使用、危険の除去、残留する危険の伝達を一体に検討する。
\begin{tcolorbox}[title=確認した法令時点]
ここは各出題原典の第1条・第2条の引用と平成30年9月の消費者庁逐条解説を比較した説明である。
e-Gov現行本文の取得は未完了で、時効等を含む出題年から現行までの全改正履歴を確認した説明ではない。
個別問題ではこの確認範囲を部分校正として追跡する。
\end{tcolorbox}

\section{単体法、双対、多目的設計}
\label{theory:dependability:optimization}
線形計画は線形の目的と制約で資源配分を表す。単体法は基底を交換し可行な端点間を移る。
退化があると単純な「毎回必ず厳密改善する」という説明は成立せず、循環を避ける選択則が必要になる。
例えば主問題を\(\max c^{\mathsf T}x,\ Ax\le b,\ x\ge0\)と置けば、
双対は\(\min b^{\mathsf T}y,\ A^{\mathsf T}y\ge c,\ y\ge0\)である。
両方の可行解について
\[
 c^{\mathsf T}x\le y^{\mathsf T}Ax\le y^{\mathsf T}b
\]
なので双対は主目的の上界を与える。可行で有限な最適値のある線形計画で両最適値は一致する。\cite{design_mit_lp}
逆問題は観測から原因を推定する別の問題設定であり、双対の呼称の代わりにはならない。
多目的最小化では、費用、重量、性能、環境負荷などの改善が相反するトレードオフが生じる。
パレート最適とは、全目的で悪化せず少なくとも一目的を厳密に改善する他の可行解が存在しないことである。
「全目的を同時に厳密改善できない」だけは弱パレートの条件である。
可行目的値\(\{(0,1),(0,0)\}\)で(0,1)は弱パレートだが(0,0)に支配されるため通常のパレートではない。

\section{信頼性・設計安全の典型的な読み方}
回路では成功経路、FTAでは故障原因を読み、同じ図形でも事象の向きを取り違えない。
積を使う前に独立性を、根を除く前に確率の範囲と非零条件を確認する。
保全では故障前後と処置時期の判断基準、法令では目的と定義、安全率では分子分母と資料時点を明示する。
過去問の組合せ表は候補の順を年度ごとに確認し、同じ概念を学ぶことと同じ候補番号を記憶することを分ける。

\section{品質を工程と判断の関係で考える}
\label{theory:quality:background}
品質管理で問うのは、ばらつきの小ささだけではない。何を満たすべきかという要求、
どのように作るかという工程、どの情報で異常を判断するかという観測を結ぶことである。
良い製品を作ること、工程が安定していること、出荷ロットを受け入れることは判断対象が違う。
目標を設定し、工程を観測して改善する流れの中で、確率統計は限られた測定から判断する道具となる。

\section{平均・分散・共分散から合成する}
\label{theory:quality:variance}
有限な二次モーメントを持つ変数を列ベクトル \(X\)、係数を \(a\) とする。
線形結合 \(Y=a^{\mathsf T}X\) の平均は \(E[Y]=a^{\mathsf T}E[X]\) である。
平均からの偏差を二乗して展開すると
\[
 \operatorname{Var}(Y)
 =\sum_i a_i^2\sigma_i^2+2\sum_{i<j}a_i a_j\operatorname{Cov}(X_i,X_j)
 =a^{\mathsf T}\Sigma a
\]
となる。\(\sigma_i^2\) は成分の分散、\(\Sigma\) は共分散行列である。
平均の線形性には独立性は不要である。分散から共分散項を消すには無相関が必要で、
独立性はその十分条件となる。独立な差 \(X-Y\) でも係数 \(-1\) が二乗されるため
分散は和となる。共通加工機の温度変動などがあれば独立とは限らず、共分散を捨てない。
単位も二乗され、長さの分散は長さの二乗、標準偏差は長さである。\cite{stat_psu_linear}

\section{正規分布の閉性と上側確率}
\label{theory:quality:normal}
\(X\sim N(\mu,\sigma^2)\) の第二引数は分散とする。
独立な正規変数の線形結合も正規となる。これは特性関数
\(\exp(i\mu t-\sigma^2t^2/2)\) の積を取れば、平均は係数付き和、
分散は係数の二乗付き和になることから分かる。
一般には各成分が正規であるだけでは足りず、同時分布の条件も必要である。
標準化 \(Z=(X-\mu)/\sigma\) により
\[
 \Pr(X\ge x)=1-\Phi\!\left(\frac{x-\mu}{\sigma}\right)
\]
と表せる。\(\Phi\) は標準正規累積分布であり、右側の面積を求める。
試験表が「上側確率」を示すか、左側の累積確率を示すかを先に読む。
差の平均に符号を付け、分散を合成し、平方根で標準偏差を出してから標準化する。\cite{stat_nist_normal}

\section{標本平均と総量は別の変数である}
\label{theory:quality:sample_mean}
同じ平均 \(\mu\)、分散 \(\sigma^2\) を持つ独立な \(n\) 個について、
総量の期待値は \(n\mu\)、分散は \(n\sigma^2\) である。
標本平均は総量を \(n\) で割るため、期待値は \(\mu\)、分散は
\(\sigma^2/n\)、標準誤差は \(\sigma/\sqrt n\) となる。
物理量の平均を取っても期待値を標本数で割って小さくするわけではない。
原変数が正規なら有限の \(n\) でもこの平均は正規である。
中心極限定理による大標本の近似とは別の根拠である。
相関のある標本では共分散項が残り、標本数を増やすだけで同じ割合で不確かさが減るとは限らない。

\section{応力・強度から故障確率へ}
\label{theory:quality:stress_strength}
強度 \(R\) と作用応力 \(S\) の単位を揃え、限界状態 \(G=R-S\) を定める。
独立正規モデルでは
\[
 b=\frac{\mu_R-\mu_S}{\sqrt{\sigma_R^2+\sigma_S^2}},\qquad
 \Pr(G<0)=\Phi(-b)
\]
となる。\(b\) は平均の余裕をばらつきで規格化した無次元の指標である。
平均強度が大きくてもばらつきが大きければ故障確率は大きくなり得る。
安全率の比とこの確率を同一視しない。
非正規の裾や強度と応力の依存があると、この式の直接適用はできない。
接続部のすき間も内径から外径を引く限界状態であり、平均の大小は全個体の符号を保証しない。

\section{どの量を数えるかで分布を選ぶ}
\label{theory:quality:models}
同じ「ばらつき」でも、測る量の範囲が違う。公平なサイコロは有限集合の離散一様である。
一定の成功確率 \(p\) の独立試行を \(n\) 回行った成功数は二項分布で、
\[
 \Pr(K=k)={n\choose k}p^k(1-p)^{n-k},\qquad
 E[K]=np,\quad \operatorname{Var}(K)=np(1-p)
\]
である。製造数を固定して不適合品を数えるモデルもこれに対応する。
コインが公平でなくても、一定 \(p\) と独立性があれば二項である。
品質の共通原因による連発や、試行ごとの確率変化があればこの条件を満たさない。
\cite{stat_nist_binomial}

\section{ポアソンの回数と指数の待ち時間}
\label{theory:quality:poisson}
二項分布で機会を増やし、小確率にして \(np=\lambda\) を有限に保つと、
個数は \(\Pr(K=k)=e^{-\lambda}\lambda^k/k!\) のポアソンへ近づく。
平均と分散はともに \(\lambda\) である。
一定率 \(r\) のポアソン過程なら時間 \(t\) の回数の平均は \(rt\) で、
一度も生じない確率が \(e^{-rt}\) となる。
したがって次の事象までの待ち時間 \(T\) は
\(\Pr(T>t)=e^{-rt}\) の指数分布となり、平均は \(1/r\) である。
率 \(r\) と区間内平均 \(\lambda=rt\) の単位を区別する。
災害や交通事故の実測に対する適合は別に検証する必要がある。
大きな \(\lambda\) ではポアソンの歪度 \(1/\sqrt{\lambda}\) が小さくなり、
平均 \(\lambda\)、分散 \(\lambda\) の正規近似が有効になり得る。
「回数だから正規近似は常に不可能」とはしない。\cite{stat_nist_poisson,stat_nist_exponential}

\section{要求品質を設計特性へ展開する}
\label{theory:quality:qfd}
利用者の言葉である要求品質を、設計で調整する品質要素へ変換するのが品質機能展開である。
試験の品質表では要求 \(i\) の重要度 \(w_i\) と要素 \(j\) との関係係数 \(c_{ij}\) を使い、
\(v_j=\sum_i w_i c_{ij}\) を求める。単なる丸の数ではなく、重い要求との対応が寄与する。
係数尺度は判断の結果に影響するため、指定された◎・○・△の値を守る。
この加重得点は実測の成功確率や安全性そのものではない。
関連する特性間の相反もあり、得点最大の列だけを設計の唯一の目標にはしない。
\cite{quality_qfd_scale}

\section{検査と上流の品質保証}
\label{theory:quality:assurance}
検査は要求への適合を判定する活動である。検査で不適合を見つけても、
それだけで不適合を作る原因は除かれない。源流管理は仕様や工程など上流にさかのぼって
発生原因を抑える考え方である。
標準化で手順と判断を明確にし、維持管理と改善を結ぶ。
工程の標準は経験的な固定ルールではなく、要求・観測・変更理由を追跡できるものとする。
QC工程表は工程と管理項目などを結ぶ計画の道具であり、統計量を時間順に並べる管理図と違う。
\cite{quality_process_jsqc}

\section{PDCAをフィードバックとして読む}
\label{theory:quality:pdca}
Planは目標と達成するプロセスを定め、Doは実施、Checkは方針・目標・要求に対する
結果を監視・測定し、Actは性能を改善する行動を取る段階である。
目標なしの測定、測定なしの改善では循環の根拠が失われる。
CをChangeやChoose、AをAdjustやAchieveという近い一般語へ置き換えない。
Actの結果を次の計画と標準へ戻すことに技術的意味がある。
\cite{quality_iso_process}

\section{2008年版を扱う問題の資料時点}
\label{theory:quality:historical}
H27の設問はISO 9001:2008/JIS Q 9001:2008を指定する。
ISOの当時の実施ガイダンスN836でもPlan-Do-Check-Actの利用が示される。
後版である2015年版の公式解説は、プロセスの管理とリスクに基づく思考を併せて説明する。
ここでは一般の四段階を両資料で照合したのであり、2008年版の全条文を取得・照合した
とは扱わない。現行規格の詳細を求める場合はその版・追補を別途確認する。
\cite{quality_iso2008_guidance,quality_iso_process}

\section{工程能力と規格限界}
\label{theory:quality:capability}
工程能力は、安定した工程が作る特性のばらつきと規格の許容幅を比較する。
正規モデルで平均 \(\mu\)、標準偏差 \(\sigma\)、規格上限 \(U\)、下限 \(L\) とすると
\[
 C_p=\frac{U-L}{6\sigma},\qquad
 C_{pk}=\min\left(\frac{U-\mu}{3\sigma},\frac{\mu-L}{3\sigma}\right)
\]
である。\(C_p\) は幅の比較だけなので平均の片寄りを直接反映しない。
\(C_{pk}\) は近い規格限界までの余裕を見る。
両者は無次元であり、標準偏差そのものではない。
非正規、工程の漂移、異なる集団の混合では、指数から不適合率への読み替えに注意する。
一定の数値を全産業共通の合格基準として勝手に設定しない。\cite{quality_nist_capability}

\section{管理図・ヒストグラム・散布図}
\label{theory:quality:charts}
管理図は工程の特性又は統計量を時間・標本順に並べ、安定工程の変動を基準に異常を探す。
中心線と管理限界は観測モデルに基づき、利用者の規格限界とは目的が違う。
限界内に全点があっても、片側への連続や系統的変化があれば安定とは限らない。
ヒストグラムは一変数の測定値を区間に分けて分布を示す。
区間幅が異なるなら高さを度数密度にし、長方形の面積を度数に対応させる。
散布図は対になった二変数の関連を示す。
分布、時間変動、二変数の関連という軸を分ければ、図の名を暗記するだけでなく用途が分かる。
\cite{quality_nist_control,quality_nist_histogram}

\section{相関は線形関連を測る}
\label{theory:quality:correlation}
有限かつ非零の分散について相関は
\(\rho=\operatorname{Cov}(X,Y)/(\sigma_X\sigma_Y)\) であり、
Cauchy--Schwarzから \(-1\le\rho\le1\) となる。
端点は完全な線形関係で達成されるため除外しない。
独立なら無相関だが、無相関だけで独立とは限らない。
対称な \(X\) と \(Y=X^2\) はその反例で、三次モーメントが零でも関数の依存が残る。
相関の絶対値は線形関連の強さを表し、因果の方向や共通原因を判別しない。
正の相関が1へ近づく説明を、任意のデータが完全な直線上にあるという断定へ変えない。

\section{回帰と決定係数の条件}
\label{theory:quality:regression}
回帰は説明変数から応答を予測するモデルを作り、残差などで適合を評価する。
切片を含む一説明変数の通常の最小二乗法では、偏差平方和を
回帰が説明する部分と残差の部分に分けられ、決定係数は \(R^2=r^2\) である。
相関の絶対値 \(|r|\) とは違う。
切片なし、多変数、異なるモデルで同じ式を機械的に使わない。
主成分分析は複数変数を合成して分散を捉える次元縮約であり、指定された応答を予測する
回帰と目的が違う。相関と回帰のどちらも観測だけから因果の証明にはならない。
\cite{stat_psu_correlation}

\section{ロットと抜取検査の二つの危険}
\label{theory:quality:acceptance}
全数を検査できない場合、サンプルからロットを受け入れるか判断する。
計数値では不適合品などを数え、計量値では長さなどの連続量を使う。
良いロットを拒否する生産者危険と、悪いロットを受け入れる消費者危険を分ける。
母集団の不適合率 \(p\) を決めた条件付き受入確率を扱うので、
「受け入れたロットが良い確率」という逆向きの条件付き確率とは違う。
\cite{quality_nist_sampling}

\section{OC曲線から閾値の効果を導く}
\label{theory:quality:oc}
大きなロットから近似的に独立に \(n\) 個を抽出し、不適合数 \(K\le c\) なら合格とする。
受入確率は二項の累積で
\[
 P_a(p)=\sum_{k=0}^c {n\choose k}p^k(1-p)^{n-k}
\]
である。良い品質点を \(p_{\mathrm g}\)、悪い品質点を \(p_{\mathrm b}\) とすれば、
\(\alpha=1-P_a(p_{\mathrm g})\)、\(\beta=P_a(p_{\mathrm b})\) となる。
\(n\) を固定して \(c\) を増やすと受入集合が広がるため、\(\alpha\) は下がり
\(\beta\) は上がる。この包含関係がトレードオフの根拠である。
\((n,c)\) の両方を固定すれば危険も固定であり、原典の「手順固定」の範囲は注意を要する。
小さい有限ロットからの非復元抽出は超幾何となり、二項の独立近似を無条件に使わない。
\cite{quality_nist_oc}

\section{ロット間の切換えと複数回の抽出}
\label{theory:quality:switching}
先行ロットの成績で後続の検査をなみ検査・きつい検査へ切り換える考え方と、
一つのロットで段階的に標本を追加する多回抜取を区別する。
後者は標本を追加したうえで合否を決める手順である。
切換えの回数や復帰条件は適用する規格に従うものであり、今回取得していない
出題時のJIS全文の数値規則を補わない。\cite{quality_nist_sampling}

\section{仮説検定は誤りを管理する判断である}
\label{theory:quality:hypothesis}
帰無仮説 \(H_0\) を基準とし、対立仮説に対してどれほど不整合な観測かを判断する。
真の \(H_0\) を棄却するのが第1種、偽の \(H_0\) を棄却しないのが第2種の誤りである。
有意水準 \(\alpha\) は第1種の誤りを管理する事前の基準であり、
観測後に \(H_0\) が真である確率ではない。
棄却されなかったことも \(H_0\) の真の証明にならない。
第2種の確率 \(\beta\) は具体的な対立側の値に依存し、検出力は \(1-\beta\) である。
同じ二種類の誤りが、品質の受入判断では生産者と消費者の危険に対応する。
\cite{quality_nist_sampling}

\section{t統計量と信頼区間の被覆}
\label{theory:quality:t_interval}
独立な正規標本で母分散が未知なら、不偏分散
\(s^2=\sum_i(X_i-\bar X)^2/(n-1)\) を使う。
平均を推定した残差は和が零という一つの制約を持つため自由度は \(n-1\) となる。
\[
 T=\frac{\bar X-\mu_0}{s/\sqrt n},\qquad
 \bar X\pm t_{1-\alpha/2,n-1}\frac{s}{\sqrt n}
\]
は、順に \(H_0\) 下での一標本t統計量と母平均の両側信頼区間である。
信頼係数は \(1-\alpha\) であって \(\alpha\) の逆数ではない。
反復標本で作る区間のうち母平均を覆う割合を表すので、観測後の一つの固定区間について
母平均がそこに存在する事後確率とはしない。\cite{stat_nist_interval}

\section{典型的な問われ方}
寸法・力の計算は変数の和差と分散の合成を先に決め、正規表で候補の範囲を判定する。
品質表は指定された関係尺度を行の重要度と掛け、列を合計する。
知識問題は要求適合、工程監視、工程能力、受入検査の判断対象を分ける。
各年度の組合せ表の列と番号は原典で照合し、同じ理論でも候補番号が同じと仮定しない。
ここで扱う例は今回確認した問題群であり、全試験回の頻度の確定は別の集計に従う。


\section{先行関係と並行作業：機械と部品の二つの順序}
\label{theory:discrete:precedence}
複数の機械で複数の部品を加工するとき、同じ機械を同時に使わない資源条件と、同じ部品の加工順の条件を別々に書く。機械AとBが部品p,q,rを各p→q→r順に処理し、各部品がA→Bへ進むなら、A列の順、B列の順、各部品のAからBへの順が全て必要である。

アロー図の矢印は作業、節点はその完了・開始の条件を表す。ある節点へ複数の矢印が入るなら、その全てが完了してから次へ進む。ダミー作業は実際の加工を増やすのでなく、必要な先行関係を追加する。例えばAqとBpの両方が終わってからBqを始めたい場合、Bpの終了節点へAqの終了からダミーを入れればよい。

候補の割当ては、見た目の左右で判断せず、各機械の順序と各部品の順序を一つずつ照合する。不要な待ちを追加するだけの場合と、必要な順を保証しない場合も異なる。並行できるという記述を全作業同時開始と考えず、共有する機械・未完了の前工程の有無を調べる。


\section{過去問と個別の解法}

\section{M/M/1は待ちと処理を分ける}\label{theory:design:queue}
到着がPoisson過程、処理時間が指数分布、窓口1、十分な待ち場所の定常モデルを考える。到着率\(\lambda\)、サービス率\(\mu\)は同じ時間単位へそろえる。平均処理時間がtなら\(\mu=1/t\)である。利用率\(\rho=\lambda/\mu<1\)が定常分布の条件である。出生死滅の釣合いから\(p_n=(1-\rho)\rho^n\)となり、
\[L=\sum_{n=0}^{\infty}np_n=\frac{\rho}{1-\rho},\qquad W=\frac L\lambda=\frac1{\mu-\lambda}.\]
待ちだけなら\(W_q=W-1/\mu\)、\(L_q=\lambda W_q=\rho^2/(1-\rho)\)。系内人数Lと行列人数Lqを混同しない。旧原典はLを待ち行列長と称するが、補助計算\(L/\mu\)はWqに等しく、処理時間を足す手順は正しい。名称と数式の検証を分ける。\cite{design48_queue_ref}
\section{PERTの最長経路と総余裕}\label{theory:design:pert}
先行作業が複数あれば全て終了してから後続を始める。結合点の最早時刻は\(T_E(j)=\max_i[T_E(i)+t_{ij}]\)、最遅時刻は終点から\(T_L(i)=\min_j[T_L(j)-t_{ij}]\)で求める。ダミー作業は所要0で先行関係を伝える。総余裕\(TF_{ij}=T_L(j)-T_E(i)-t_{ij}\)が0の作業はクリティカルである。
全工期は最長経路長であり、最も安い作業を無条件に縮めても工期は下がらない。複数の最長経路があれば全てを短くする。平成30年度はa3とb2が同じ22日で二つの最長枝を作るため、その作業和集合が管理対象になる。短縮前後の最長経路を再計算し、別経路への支配交代も確かめる。
\section{線形計画の実行可能性と頂点}\label{theory:design:lp}
生産量を変数にし、資源上限を不等式、利益を目的関数へ置き換える。利益係数が最大の製品だけを増やす前に、全材料と時間の制約を確認する。二変数なら制約境界の交点を求め、実行可能領域の頂点の利益を比較する。台数が整数なら連続最適解をそのまま使わず整数候補を確認する。
令和2年度は\(3x+2y\le24\)、\(x+3y\le15\)の交点\((6,3)\)が当初の最適である。利益係数変更後もこの頂点を最適にするには、隣接頂点\((0,5),(8,0)\)と比較して
\[21+6\Delta c\ge15,\qquad21+6\Delta c\ge16+8\Delta c,\qquad-1\le\Delta c\le5/2.\]
端点では最適解が複数になる。元の点が最適であることと唯一であることを区別する。全組合せ最適化に指数時間下界が証明済みと主張せず、解ける部分クラスと近似品質の条件を別に扱う。
\section{ユニバーサルデザインの利用者と場面}\label{theory:design:universal}
UDは可能な限り多くの人が特別な変更なしに使える製品・環境を目指す。1997年作業班の七原則は、公平な利用、柔軟な利用、単純で直観的な利用、認知できる情報、失敗への寛容、少ない身体的努力、接近・利用に必要な大きさと空間である。原則名を対象例へ対応させ、自動ドアの省力、速度変更の柔軟性、多形式情報、低い位置からの到達を比較する。安全のため特定利用者を一律に排除する機能を、そのままUDの例にはしない。\cite{design48_ud}
\section{分割と機能独立性}\label{theory:design:principles}
多様な要求の調和を図り、複雑系を適切な機能単位へ分割して管理する。複雑であるほど優れた設計という結論は出ない。機能間の結合を減らすことは変更影響を局所化するが、インタフェース・全体相互作用を無視できるわけではない。ばらつきは設計対象であり理想値だけを実機の保証としない。
\section{角法と図面情報}\label{theory:design:drawing}
第三角では見る側と同じ側へ図を置き、平面図は上、右側面図は右。第一角では反対側となり平面図は下、左側面図は右になる。上と下の一語が異なる候補を全体の印象で正にしない。ISOは国際標準化機構で国別のイタリア規格ではない。公開規格には第一角・第三角の節があるが有料節の全図まで読了したとはしない。\cite{design48_iso}
\section{要求性能から構造詳細と照査へ}\label{theory:design:performance}
合理的構造体を実現するため、要求性能を設定し、構造計画、構造詳細の設定、性能照査へ進む。構造詳細は部材寸法・材料等の照査に必要な情報であり、詳細設計という一般工程名と異なる。供用期間の作用と劣化を踏まえて、設定した要求が満たされるかを評価する。使用性能という一つの性能だけに全要求を限定しない。\cite{design48_mlit}

\section{利得表は他者の行動ごとに比較する}
\label{theory:design:game}
各参加者が自分の利得を最大にするゲームでは、相手の行動を固定して自分の候補を比較する。相手のどの行動でも優る戦略が支配戦略である。両者の最適反応が同時に満たされる組合せは、片方だけ変更しても利得を増せないナッシュ均衡となる。行プレイヤーの利得を列ごと、列プレイヤーの利得を行ごとに比較し、表を転置して読む誤りを避ける。両者の合計利得の最大化は協調の評価で、独立の最適反応とは別である。

\section{期待値と決定ノードの後向き評価}
\label{theory:design:decision}
確率が既知の状態\(s\)について期待便益は\(\sum_s p_s B_s\)。一定費用\(C\)を引けば期待純便益になる。決定後に状態を観測して対応を選べる木では、末端から機会ノードを期待値、決定ノードを最大便益又は最小損失で置き換える。選べる二案を半々に平均してはならない。先に選ぶ案と、状態を見てから選べる案では利用できる情報が違う。

\section{限界利益と検査の期待費用}
\label{theory:design:economics}
一人又は一個当たり売上を\(p\)、変動費を\(v\)、期間固定費を\(F\)、目標利益を\(G\)とすると、必要数量は\((F+G)/(p-v)\)。\(p-v>0\)が前提である。最後に営業日数・設備数で割って、求める単位へ変換する。
\label{theory:design:inspection}
検査回数\(x\)で不具合確率が\(p(x)\)、一回費用\(c\)、損害額\(D\)なら期待総費用は\(cx+Dp(x)\)。検査を増やす費用と確率低下の利益を比較する。連続最適点を微分で求めた後、実際に許される整数値と端点を調べる。推定確率のモデルをそのまま現場の保証値とはしない。

\section{凸性は最適解の意味を強める}
\label{theory:design:convex}
線形計画は目的と制約がともに線形である。凸集合上で凸関数を最小化する凸計画では局所最適点が大域最適点となる\autocite{finish56_convex}。目的だけ凸でも実行可能集合が非凸ならその性質は保証されない。整数制約は実行可能集合を離散化し、連続問題の解を丸めるだけでは制約を満たさない場合がある。ヒューリスティックは候補解を得る手段で、最適性証明や誤差保証は別に扱う。

\section{経済的発注量は二つの費用を釣り合わせる}
\label{theory:design:eoq}
一定需要\(d\)、即時一括補充、欠品なしのモデルでは、発注量\(q\)に対し発注回数\(d/q\)、平均在庫\(q/2\)。一回費用\(k\)、年間単位在庫費\(h\)なら
\begin{equation}C(q)=\frac{kd}{q}+\frac{hq}{2},\qquad q_*=\sqrt{\frac{2kd}{h}}.\end{equation}
\(q>0\)で\(C''=2kd/q^3>0\)だから停留点は最小である\autocite{finish56_eoq}。上限、ロット単位、需要変動、納期、数量割引があれば再計算する。全需要に比例して一定の購入費は、このモデルでは最適発注量に影響しない。

\section{三つの解析公式の成立条件}
\label{theory:design:mathformulas}
Taylor展開の中心が0ならMaclaurin展開で\(f(x)=\sum_{k\ge0}f^{(k)}(0)x^k/k!\)。無限回微分可能だけでは級数と元の関数の一致は保証されず、剰余の収束を含む条件が必要である。複素指数の\(e^{i\theta}=\cos\theta+i\sin\theta\)から\(\theta=\pi\)で\(e^{i\pi}+1=0\)となる。
ロピタルの定理は\(0/0\)又は適切な無限大比等の不定形で、微分可能性、分母微分非零、微分比の極限等を確認して使う。任意の比の分子・分母を微分する規則ではない。

\section{不等式の向きと等号条件}
\label{theory:design:inequality}
非負の数の幾何平均は算術平均以下で、全数が同じなら等号を取る。\([0,\pi/2]\)上の\(\sin\theta\)の凹性から弦との比較で\(\sin\theta\ge2\theta/\pi\)。\(0<\theta\le\pi/2\)なら\(\sin\theta/\theta\ge2/\pi\)、右端で等号である。
一方、区間上で\(f''>0\)なら厳密凸で、異なる二点に対し\(f((x_1+x_2)/2)<(f(x_1)+f(x_2))/2\)。端点を含む不等式と厳密凸の厳密不等号を読み分ける。

\section{歯車の速度比とトルク}
\label{theory:design:gear}
回転動力は\(P=T\omega\)。速度比\(r=\omega_{\mathrm{in}}/\omega_{\mathrm{out}}\)、効率\(\eta\)なら\(T_{\mathrm{out}}=\eta rT_{\mathrm{in}}\)。減速で角速度が下がると、同じ動力に対応するトルクは増す。比の定義を入力/出力かその逆かで確認し、回転方向と大きさの比較を分ける。

\section{平面応力と平面ひずみは異なる近似}
\label{theory:design:plane}
薄板の自由表面に対する平面応力では\(\sigma_z,\tau_{xz},\tau_{yz}\simeq0\)で、面外ひずみは一般に0でない。軸方向変形を拘束した一様な長い構造の内部に対する平面ひずみでは\(\varepsilon_z,\gamma_{xz},\gamma_{yz}\simeq0\)で、軸応力は一般に残る。等方線形弾性なら後者の\(\sigma_z=\nu(\sigma_x+\sigma_y)\)。端部、軸方向の自由伸び、異方性を含む場合は境界条件からモデルを選び直す。

\subsection{平成23年度 問1-1}
\label{q:h23-reg-basic-1-1}

\noindent\textbf{原典：} sources/H23/basic.pdf、PDF物理1頁。

\textbf{公式正答：} 2（single）。

\par\textit{原典本文の校正と独自解説は未作成。分類のみの登録項目。}

\subsection{平成23年度 問1-2}
\label{q:h23-reg-basic-1-2}

\noindent\textbf{原典：} sources/H23/basic.pdf、PDF物理1頁。

\textbf{公式正答：} 4（single）。

\par\textit{原典本文の校正と独自解説は未作成。分類のみの登録項目。}

\subsection{平成23年度 問1-3}
\label{q:h23-reg-basic-1-3}

\noindent\textbf{原典：} sources/H23/basic.pdf、PDF物理2頁。

\textbf{公式正答：} 3（single）。

\par\textit{原典本文の校正と独自解説は未作成。分類のみの登録項目。}

\subsection{平成23年度 問1-4}
\label{q:h23-reg-basic-1-4}

\noindent\textbf{原典：} sources/H23/basic.pdf、PDF物理3頁。

\textbf{公式正答：} 1（single）。

\par\textit{原典本文の校正と独自解説は未作成。分類のみの登録項目。}

\subsection{平成23年度 問1-5}
\label{q:h23-reg-basic-1-5}

\noindent\textbf{原典：} sources/H23/basic.pdf、PDF物理4頁。

\textbf{公式正答：} 4（single）。

\par\textit{原典本文の校正と独自解説は未作成。分類のみの登録項目。}

\subsection{平成24年度 問1-1}
\label{q:h24-reg-basic-1-1}

\noindent\textbf{原典：} sources/H24/basic.pdf、PDF物理1頁。

\textbf{公式正答：} 4（single）。

\par\textit{原典本文の校正と独自解説は未作成。分類のみの登録項目。}

\subsection{平成24年度 問1-2}
\label{q:h24-reg-basic-1-2}

\noindent\textbf{原典：} sources/H24/basic.pdf、PDF物理2頁。

\textbf{公式正答：} 5（single）。

\par\textit{原典本文の校正と独自解説は未作成。分類のみの登録項目。}

\subsection{平成24年度 問1-3}
\label{q:h24-reg-basic-1-3}

\noindent\textbf{原典：} sources/H24/basic.pdf、PDF物理3頁。

\textbf{公式正答：} 3（single）。

\par\textit{原典本文の校正と独自解説は未作成。分類のみの登録項目。}

\subsection{平成24年度 問1-4}
\label{q:h24-reg-basic-1-4}

\noindent\textbf{原典：} sources/H24/basic.pdf、PDF物理4頁。

\textbf{公式正答：} 5（single）。

\par\textit{原典本文の校正と独自解説は未作成。分類のみの登録項目。}

\subsection{平成24年度 問1-5}
\label{q:h24-reg-basic-1-5}

\noindent\textbf{原典：} sources/H24/basic.pdf、PDF物理4頁。

\textbf{公式正答：} 5（single）。

\par\textit{原典本文の校正と独自解説は未作成。分類のみの登録項目。}

\subsection{平成25年度 問1-1}
\label{q:h25-reg-basic-1-1}

\noindent\textbf{原典：} sources/H25/basic.pdf、PDF物理1頁。

\textbf{公式正答：} 3（single）。

\par\textit{原典本文の校正と独自解説は未作成。分類のみの登録項目。}

\subsection{平成25年度 問1-2}
\label{q:h25-reg-basic-1-2}

\noindent\textbf{原典：} sources/H25/basic.pdf、PDF物理2頁。

\textbf{公式正答：} 5（single）。

\par\textit{原典本文の校正と独自解説は未作成。分類のみの登録項目。}

\subsection{平成25年度 問1-3}
\label{q:h25-reg-basic-1-3}

\noindent\textbf{原典：} sources/H25/basic.pdf、PDF物理3頁。

\textbf{公式正答：} 2（single）。

\par\textit{原典本文の校正と独自解説は未作成。分類のみの登録項目。}

\subsection{平成25年度 問1-4}
\label{q:h25-reg-basic-1-4}

\noindent\textbf{原典：} sources/H25/basic.pdf、PDF物理4頁。

\textbf{公式正答：} 4（single）。

\par\textit{原典本文の校正と独自解説は未作成。分類のみの登録項目。}

\subsection{平成25年度 問1-5}
\label{q:h25-reg-basic-1-5}

\noindent\textbf{原典：} sources/H25/basic.pdf、PDF物理5頁。

\textbf{公式正答：} 4（single）。

\par\textit{原典本文の校正と独自解説は未作成。分類のみの登録項目。}

\subsection{平成25年度 問1-6}
\label{q:h25-reg-basic-1-6}

\noindent\textbf{原典：} sources/H25/basic.pdf、PDF物理5頁。

\textbf{公式正答：} 3（single）。

\par\textit{原典本文の校正と独自解説は未作成。分類のみの登録項目。}

\subsection{平成26年度 問1-1}
\label{q:h26-reg-basic-1-1}

\noindent\textbf{原典：} sources/H26/basic.pdf、PDF物理1頁。

\textbf{公式正答：} 5（single）。

\par\textit{原典本文の校正と独自解説は未作成。分類のみの登録項目。}

\subsection{平成26年度 問1-2}
\label{q:h26-reg-basic-1-2}

\noindent\textbf{原典：} sources/H26/basic.pdf、PDF物理1頁。

\textbf{公式正答：} 4（single）。

\par\textit{原典本文の校正と独自解説は未作成。分類のみの登録項目。}

\subsection{平成26年度 問1-3}
\label{q:h26-reg-basic-1-3}

\noindent\textbf{原典：} sources/H26/basic.pdf、PDF物理2頁。

\textbf{公式正答：} 1（single）。

\par\textit{原典本文の校正と独自解説は未作成。分類のみの登録項目。}

\subsection{平成26年度 問1-4}
\label{q:h26-reg-basic-1-4}

\noindent\textbf{原典：} sources/H26/basic.pdf、PDF物理3頁。

\textbf{公式正答：} 3（single）。

\par\textit{原典本文の校正と独自解説は未作成。分類のみの登録項目。}

\subsection{平成26年度 問1-5}
\label{q:h26-reg-basic-1-5}

\noindent\textbf{原典：} sources/H26/basic.pdf、PDF物理3頁。

\textbf{公式正答：} 1（single）。

\par\textit{原典本文の校正と独自解説は未作成。分類のみの登録項目。}

\subsection{平成26年度 問1-6}
\label{q:h26-reg-basic-1-6}

\noindent\textbf{原典：} sources/H26/basic.pdf、PDF物理4頁。

\textbf{公式正答：} 3（single）。

\par\textit{原典本文の校正と独自解説は未作成。分類のみの登録項目。}

\subsection{平成27年度 問1-1}
\label{q:h27-reg-basic-1-1}

\noindent\textbf{原典：} sources/H27/basic.pdf、PDF物理1頁。

\textbf{公式正答：} 5（single）。

\par\textit{原典本文の校正と独自解説は未作成。分類のみの登録項目。}

\subsection{平成27年度 問1-2}
\label{q:h27-reg-basic-1-2}

\noindent\textbf{原典：} sources/H27/basic.pdf、PDF物理2頁。

\textbf{公式正答：} 4（single）。

\par\textit{原典本文の校正と独自解説は未作成。分類のみの登録項目。}

\subsection{平成27年度 問1-3}
\label{q:h27-reg-basic-1-3}

\noindent\textbf{原典：} sources/H27/basic.pdf、PDF物理2頁。

\textbf{公式正答：} 5（single）。

\par\textit{原典本文の校正と独自解説は未作成。分類のみの登録項目。}

\subsection{平成27年度 問1-4}
\label{q:h27-reg-basic-1-4}

\noindent\textbf{原典：} sources/H27/basic.pdf、PDF物理3頁。

\textbf{公式正答：} 3（single）。

\par\textit{原典本文の校正と独自解説は未作成。分類のみの登録項目。}

\subsection{平成27年度 問1-5}
\label{q:h27-reg-basic-1-5}

\noindent\textbf{原典：} sources/H27/basic.pdf、PDF物理4頁。

\textbf{公式正答：} 1（single）。

\par\textit{原典本文の校正と独自解説は未作成。分類のみの登録項目。}

\subsection{平成27年度 問1-6}
\label{q:h27-reg-basic-1-6}

\noindent\textbf{原典：} sources/H27/basic.pdf、PDF物理4頁。

\textbf{公式正答：} 5（single）。

\par\textit{原典本文の校正と独自解説は未作成。分類のみの登録項目。}

\subsection{平成28年度 問1-1}
\label{q:h28-reg-basic-1-1}

\noindent\textbf{原典：} sources/H28/basic.pdf、PDF物理1頁。

\textbf{公式正答：} 4（single）。

\par\textit{原典本文の校正と独自解説は未作成。分類のみの登録項目。}

\subsection{平成28年度 問1-2}
\label{q:h28-reg-basic-1-2}

\noindent\textbf{原典：} sources/H28/basic.pdf、PDF物理2頁。

\textbf{公式正答：} 4（single）。

\par\textit{原典本文の校正と独自解説は未作成。分類のみの登録項目。}

\subsection{平成28年度 問1-3}
\label{q:h28-reg-basic-1-3}

\noindent\textbf{原典：} sources/H28/basic.pdf、PDF物理3頁。

\textbf{公式正答：} 4（single）。

\par\textit{原典本文の校正と独自解説は未作成。分類のみの登録項目。}

\subsection{平成28年度 問1-4}
\label{q:h28-reg-basic-1-4}

\noindent\textbf{原典：} sources/H28/basic.pdf、PDF物理4頁。

\textbf{公式正答：} 4（single）。

\par\textit{原典本文の校正と独自解説は未作成。分類のみの登録項目。}

\subsection{平成28年度 問1-5}
\label{q:h28-reg-basic-1-5}

\noindent\textbf{原典：} sources/H28/basic.pdf、PDF物理5頁。

\textbf{公式正答：} 5（single）。

\par\textit{原典本文の校正と独自解説は未作成。分類のみの登録項目。}

\subsection{平成28年度 問1-6}
\label{q:h28-reg-basic-1-6}

\noindent\textbf{原典：} sources/H28/basic.pdf、PDF物理6頁。

\textbf{公式正答：} 1（single）。

\par\textit{原典本文の校正と独自解説は未作成。分類のみの登録項目。}

\subsection{平成29年度 問1-1}
\label{q:h29-reg-basic-1-1}

\noindent\textbf{原典：} sources/H29/basic.pdf、PDF物理1頁。

\textbf{公式正答：} 5（single）。

\par\textit{原典本文の校正と独自解説は未作成。分類のみの登録項目。}

\subsection{平成29年度 問1-2}
\label{q:h29-reg-basic-1-2}

\noindent\textbf{原典：} sources/H29/basic.pdf、PDF物理1頁。

\textbf{公式正答：} 5（single）。

\par\textit{原典本文の校正と独自解説は未作成。分類のみの登録項目。}

\subsection{平成29年度 問1-3}
\label{q:h29-reg-basic-1-3}

\noindent\textbf{原典：} sources/H29/basic.pdf、PDF物理2頁。

\textbf{公式正答：} 2（single）。

\par\textit{原典本文の校正と独自解説は未作成。分類のみの登録項目。}

\subsection{平成29年度 問1-4}
\label{q:h29-reg-basic-1-4}

\noindent\textbf{原典：} sources/H29/basic.pdf、PDF物理3頁。

\textbf{公式正答：} 5（single）。

\par\textit{原典本文の校正と独自解説は未作成。分類のみの登録項目。}

\subsection{平成29年度 問1-5}
\label{q:h29-reg-basic-1-5}

\noindent\textbf{原典：} sources/H29/basic.pdf、PDF物理4頁。

\textbf{公式正答：} 1（single）。

\par\textit{原典本文の校正と独自解説は未作成。分類のみの登録項目。}

\subsection{平成29年度 問1-6}
\label{q:h29-reg-basic-1-6}

\noindent\textbf{原典：} sources/H29/basic.pdf、PDF物理5頁。

\textbf{公式正答：} 2（single）。

\par\textit{原典本文の校正と独自解説は未作成。分類のみの登録項目。}

\subsection{平成30年度 問1-1}
\label{q:h30-reg-basic-1-1}

\noindent\textbf{原典：} sources/H30/basic.pdf、PDF物理1頁。

\textbf{公式正答：} 3（single）。

\par\textit{原典本文の校正と独自解説は未作成。分類のみの登録項目。}

\subsection{平成30年度 問1-2}
\label{q:h30-reg-basic-1-2}

\noindent\textbf{原典：} sources/H30/basic.pdf、PDF物理2頁。

\textbf{公式正答：} 1（single）。

\par\textit{原典本文の校正と独自解説は未作成。分類のみの登録項目。}

\subsection{平成30年度 問1-3}
\label{q:h30-reg-basic-1-3}

\noindent\textbf{原典：} sources/H30/basic.pdf、PDF物理3頁。

\textbf{公式正答：} 2（single）。

\par\textit{原典本文の校正と独自解説は未作成。分類のみの登録項目。}

\subsection{平成30年度 問1-4}
\label{q:h30-reg-basic-1-4}

\noindent\textbf{原典：} sources/H30/basic.pdf、PDF物理4頁。

\textbf{公式正答：} 2（single）。

\par\textit{原典本文の校正と独自解説は未作成。分類のみの登録項目。}

\subsection{平成30年度 問1-5}
\label{q:h30-reg-basic-1-5}

\noindent\textbf{原典：} sources/H30/basic.pdf、PDF物理5頁。

\textbf{公式正答：} 3（single）。

\par\textit{原典本文の校正と独自解説は未作成。分類のみの登録項目。}

\subsection{平成30年度 問1-6}
\label{q:h30-reg-basic-1-6}

\noindent\textbf{原典：} sources/H30/basic.pdf、PDF物理5頁。

\textbf{公式正答：} 4（single）。

\par\textit{原典本文の校正と独自解説は未作成。分類のみの登録項目。}

\subsection{令和元年度 問1-1}
\label{q:r01-reg-basic-1-1}

\noindent\textbf{原典：} sources/R01/basic.pdf、PDF物理1頁。

\textbf{公式正答：} 5（single）。

\par\textit{原典本文の校正と独自解説は未作成。分類のみの登録項目。}

\subsection{令和元年度 問1-2}
\label{q:r01-reg-basic-1-2}

\noindent\textbf{原典：} sources/R01/basic.pdf、PDF物理2頁。

\textbf{公式正答：} 2（single）。

\par\textit{原典本文の校正と独自解説は未作成。分類のみの登録項目。}

\subsection{令和元年度 問1-3}
\label{q:r01-reg-basic-1-3}

\noindent\textbf{原典：} sources/R01/basic.pdf、PDF物理3頁。

\textbf{公式正答：} 5（single）。

\par\textit{原典本文の校正と独自解説は未作成。分類のみの登録項目。}

\subsection{令和元年度 問1-4}
\label{q:r01-reg-basic-1-4}

\noindent\textbf{原典：} sources/R01/basic.pdf、PDF物理4頁。

\textbf{公式正答：} 5（single）。

\par\textit{原典本文の校正と独自解説は未作成。分類のみの登録項目。}

\subsection{令和元年度 問1-5}
\label{q:r01-reg-basic-1-5}

\noindent\textbf{原典：} sources/R01/basic.pdf、PDF物理5頁。

\textbf{公式正答：} 2（single）。

\par\textit{原典本文の校正と独自解説は未作成。分類のみの登録項目。}

\subsection{令和元年度 問1-6}
\label{q:r01-reg-basic-1-6}

\noindent\textbf{原典：} sources/R01/basic.pdf、PDF物理6頁。

\textbf{公式正答：} 3（single）。

\par\textit{原典本文の校正と独自解説は未作成。分類のみの登録項目。}

\subsection{令和元年度再試験 問1-1}
\label{q:r01-ret-basic-1-1}

\noindent\textbf{原典：} sources/R01\_retest/basic.pdf、PDF物理1頁。

\textbf{公式正答：} 5（single）。

\par\textit{原典本文の校正と独自解説は未作成。分類のみの登録項目。}

\subsection{令和元年度再試験 問1-2}
\label{q:r01-ret-basic-1-2}

\noindent\textbf{原典：} sources/R01\_retest/basic.pdf、PDF物理2頁。

\textbf{公式正答：} 3（single）。

\par\textit{原典本文の校正と独自解説は未作成。分類のみの登録項目。}

\subsection{令和元年度再試験 問1-3}
\label{q:r01-ret-basic-1-3}

\noindent\textbf{原典：} sources/R01\_retest/basic.pdf、PDF物理3頁。

\textbf{公式正答：} 1（single）。

\par\textit{原典本文の校正と独自解説は未作成。分類のみの登録項目。}

\subsection{令和元年度再試験 問1-4}
\label{q:r01-ret-basic-1-4}

\noindent\textbf{原典：} sources/R01\_retest/basic.pdf、PDF物理4頁。

\textbf{公式正答：} 2（single）。

\par\textit{原典本文の校正と独自解説は未作成。分類のみの登録項目。}

\subsection{令和元年度再試験 問1-5}
\label{q:r01-ret-basic-1-5}

\noindent\textbf{原典：} sources/R01\_retest/basic.pdf、PDF物理5頁。

\textbf{公式正答：} 2（single）。

\par\textit{原典本文の校正と独自解説は未作成。分類のみの登録項目。}

\subsection{令和元年度再試験 問1-6}
\label{q:r01-ret-basic-1-6}

\noindent\textbf{原典：} sources/R01\_retest/basic.pdf、PDF物理5頁。

\textbf{公式正答：} 5（single）。

\par\textit{原典本文の校正と独自解説は未作成。分類のみの登録項目。}

\subsection{令和2年度 問1-1}
\label{q:r02-reg-basic-1-1}

\noindent\textbf{原典：} sources/R02/basic.pdf、PDF物理1頁。

\textbf{公式正答：} 3（single）。

\par\textit{原典本文の校正と独自解説は未作成。分類のみの登録項目。}

\subsection{令和2年度 問1-2}
\label{q:r02-reg-basic-1-2}

\noindent\textbf{原典：} sources/R02/basic.pdf、PDF物理2頁。

\textbf{公式正答：} 4（single）。

\par\textit{原典本文の校正と独自解説は未作成。分類のみの登録項目。}

\subsection{令和2年度 問1-3}
\label{q:r02-reg-basic-1-3}

\noindent\textbf{原典：} sources/R02/basic.pdf、PDF物理3頁。

\textbf{公式正答：} 3（single）。

\par\textit{原典本文の校正と独自解説は未作成。分類のみの登録項目。}

\subsection{令和2年度 問1-4}
\label{q:r02-reg-basic-1-4}

\noindent\textbf{原典：} sources/R02/basic.pdf、PDF物理4頁。

\textbf{公式正答：} 5（single）。

\par\textit{原典本文の校正と独自解説は未作成。分類のみの登録項目。}

\subsection{令和2年度 問1-5}
\label{q:r02-reg-basic-1-5}

\noindent\textbf{原典：} sources/R02/basic.pdf、PDF物理5頁。

\textbf{公式正答：} 5（single）。

\par\textit{原典本文の校正と独自解説は未作成。分類のみの登録項目。}

\subsection{令和2年度 問1-6}
\label{q:r02-reg-basic-1-6}

\noindent\textbf{原典：} sources/R02/basic.pdf、PDF物理6頁。

\textbf{公式正答：} 3（single）。

\par\textit{原典本文の校正と独自解説は未作成。分類のみの登録項目。}

\subsection{令和3年度 問1-1}
\label{q:r03-reg-basic-1-1}

\noindent\textbf{原典：} sources/R03/basic.pdf、PDF物理1頁。

\textbf{公式正答：} 4（single）。

\par\textit{原典本文の校正と独自解説は未作成。分類のみの登録項目。}

\subsection{令和3年度 問1-2}
\label{q:r03-reg-basic-1-2}

\noindent\textbf{原典：} sources/R03/basic.pdf、PDF物理2頁。

\textbf{公式正答：} 2（single）。

\par\textit{原典本文の校正と独自解説は未作成。分類のみの登録項目。}

\subsection{令和3年度 問1-3}
\label{q:r03-reg-basic-1-3}

\noindent\textbf{原典：} sources/R03/basic.pdf、PDF物理3頁。

\textbf{公式正答：} 2（single）。

\par\textit{原典本文の校正と独自解説は未作成。分類のみの登録項目。}

\subsection{令和3年度 問1-4}
\label{q:r03-reg-basic-1-4}

\noindent\textbf{原典：} sources/R03/basic.pdf、PDF物理3頁。

\textbf{公式正答：} 3（single）。

\par\textit{原典本文の校正と独自解説は未作成。分類のみの登録項目。}

\subsection{令和3年度 問1-5}
\label{q:r03-reg-basic-1-5}

\noindent\textbf{原典：} sources/R03/basic.pdf、PDF物理4頁。

\textbf{公式正答：} 5（single）。

\par\textit{原典本文の校正と独自解説は未作成。分類のみの登録項目。}

\subsection{令和3年度 問1-6}
\label{q:r03-reg-basic-1-6}

\noindent\textbf{原典：} sources/R03/basic.pdf、PDF物理5頁。

\textbf{公式正答：} 3（single）。

\par\textit{原典本文の校正と独自解説は未作成。分類のみの登録項目。}

\subsection{令和4年度 問1-1}
\label{q:r04-reg-basic-1-1}

\noindent\textbf{原典：} sources/R04/basic.pdf、PDF物理1頁。

\textbf{公式正答：} 1（single）。

\par\textit{原典本文の校正と独自解説は未作成。分類のみの登録項目。}

\subsection{令和4年度 問1-2}
\label{q:r04-reg-basic-1-2}

\noindent\textbf{原典：} sources/R04/basic.pdf、PDF物理2頁。

\textbf{公式正答：} 4（single）。

\par\textit{原典本文の校正と独自解説は未作成。分類のみの登録項目。}

\subsection{令和4年度 問1-3}
\label{q:r04-reg-basic-1-3}

\noindent\textbf{原典：} sources/R04/basic.pdf、PDF物理2頁。

\textbf{公式正答：} 3（single）。

\par\textit{原典本文の校正と独自解説は未作成。分類のみの登録項目。}

\subsection{令和4年度 問1-4}
\label{q:r04-reg-basic-1-4}

\noindent\textbf{原典：} sources/R04/basic.pdf、PDF物理3頁。

\textbf{公式正答：} 1（single）。

\par\textit{原典本文の校正と独自解説は未作成。分類のみの登録項目。}

\subsection{令和4年度 問1-5}
\label{q:r04-reg-basic-1-5}

\noindent\textbf{原典：} sources/R04/basic.pdf、PDF物理3頁。

\textbf{公式正答：} 3（single）。

\par\textit{原典本文の校正と独自解説は未作成。分類のみの登録項目。}

\subsection{令和4年度 問1-6}
\label{q:r04-reg-basic-1-6}

\noindent\textbf{原典：} sources/R04/basic.pdf、PDF物理4頁。

\textbf{公式正答：} 2（single）。

\par\textit{原典本文の校正と独自解説は未作成。分類のみの登録項目。}

\subsection{令和5年度 問1-1}
\label{q:r05-reg-basic-1-1}

\noindent\textbf{原典：} sources/R05/basic.pdf、PDF物理1頁。

\textbf{公式正答：} 5（single）。

\par\textit{原典本文の校正と独自解説は未作成。分類のみの登録項目。}

\subsection{令和5年度 問1-2}
\label{q:r05-reg-basic-1-2}

\noindent\textbf{原典：} sources/R05/basic.pdf、PDF物理2頁。

\textbf{公式正答：} 4（single）。

\par\textit{原典本文の校正と独自解説は未作成。分類のみの登録項目。}

\subsection{令和5年度 問1-3}
\label{q:r05-reg-basic-1-3}

\noindent\textbf{原典：} sources/R05/basic.pdf、PDF物理3頁。

\textbf{公式正答：} 3（single）。

\par\textit{原典本文の校正と独自解説は未作成。分類のみの登録項目。}

\subsection{令和5年度 問1-4}
\label{q:r05-reg-basic-1-4}

\noindent\textbf{原典：} sources/R05/basic.pdf、PDF物理3頁。

\textbf{公式正答：} 4（single）。

\par\textit{原典本文の校正と独自解説は未作成。分類のみの登録項目。}

\subsection{令和5年度 問1-5}
\label{q:r05-reg-basic-1-5}

\noindent\textbf{原典：} sources/R05/basic.pdf、PDF物理4頁。

\textbf{公式正答：} 1（single）。

\par\textit{原典本文の校正と独自解説は未作成。分類のみの登録項目。}

\subsection{令和5年度 問1-6}
\label{q:r05-reg-basic-1-6}

\noindent\textbf{原典：} sources/R05/basic.pdf、PDF物理4頁。

\textbf{公式正答：} 2（single）。

\par\textit{原典本文の校正と独自解説は未作成。分類のみの登録項目。}

\subsection{令和6年度 問1-1}
\label{q:r06-reg-basic-1-1}

\noindent\textbf{原典：} sources/R06/basic.pdf、PDF物理1頁。

\textbf{公式正答：} 1（single）。

\par\textit{原典本文の校正と独自解説は未作成。分類のみの登録項目。}

\subsection{令和6年度 問1-2}
\label{q:r06-reg-basic-1-2}

\noindent\textbf{原典：} sources/R06/basic.pdf、PDF物理2頁。

\textbf{公式正答：} 2（single）。

\par\textit{原典本文の校正と独自解説は未作成。分類のみの登録項目。}

\subsection{令和6年度 問1-3}
\label{q:r06-reg-basic-1-3}

\noindent\textbf{原典：} sources/R06/basic.pdf、PDF物理2頁。

\textbf{公式正答：} 1（single）。

\par\textit{原典本文の校正と独自解説は未作成。分類のみの登録項目。}

\subsection{令和6年度 問1-4}
\label{q:r06-reg-basic-1-4}

\noindent\textbf{原典：} sources/R06/basic.pdf、PDF物理3頁。

\textbf{公式正答：} 3（single）。

\par\textit{原典本文の校正と独自解説は未作成。分類のみの登録項目。}

\subsection{令和6年度 問1-5}
\label{q:r06-reg-basic-1-5}

\noindent\textbf{原典：} sources/R06/basic.pdf、PDF物理4頁。

\textbf{公式正答：} 4（single）。

\par\textit{原典本文の校正と独自解説は未作成。分類のみの登録項目。}

\subsection{令和6年度 問1-6}
\label{q:r06-reg-basic-1-6}

\noindent\textbf{原典：} sources/R06/basic.pdf、PDF物理5頁。

\textbf{公式正答：} 3（single）。

\par\textit{原典本文の校正と独自解説は未作成。分類のみの登録項目。}

\subsection{令和7年度 問1-1}
\label{q:r07-reg-basic-1-1}

\noindent\textbf{原典：} sources/R07/basic.pdf、PDF物理1頁。

\textbf{公式正答：} 3（single）。

\par\textit{原典本文の校正と独自解説は未作成。分類のみの登録項目。}

\subsection{令和7年度 問1-2}
\label{q:r07-reg-basic-1-2}

\noindent\textbf{原典：} sources/R07/basic.pdf、PDF物理2頁。

\textbf{公式正答：} 4（single）。

\par\textit{原典本文の校正と独自解説は未作成。分類のみの登録項目。}

\subsection{令和7年度 問1-3}
\label{q:r07-reg-basic-1-3}

\noindent\textbf{原典：} sources/R07/basic.pdf、PDF物理2頁。

\textbf{公式正答：} 2（single）。

\par\textit{原典本文の校正と独自解説は未作成。分類のみの登録項目。}

\subsection{令和7年度 問1-4}
\label{q:r07-reg-basic-1-4}

\noindent\textbf{原典：} sources/R07/basic.pdf、PDF物理3頁。

\textbf{公式正答：} 1（single）。

\par\textit{原典本文の校正と独自解説は未作成。分類のみの登録項目。}

\subsection{令和7年度 問1-5}
\label{q:r07-reg-basic-1-5}

\noindent\textbf{原典：} sources/R07/basic.pdf、PDF物理3頁。

\textbf{公式正答：} 5（single）。

\par\textit{原典本文の校正と独自解説は未作成。分類のみの登録項目。}

\subsection{令和7年度 問1-6}
\label{q:r07-reg-basic-1-6}

\noindent\textbf{原典：} sources/R07/basic.pdf、PDF物理4頁。

\textbf{公式正答：} 3（single）。

\par\textit{原典本文の校正と独自解説は未作成。分類のみの登録項目。}

